%%
%% CSD Assignment 1 - Final Technical Report
%% Based on the ACM "acmart" manuscript (single-column) style.
%%
\documentclass[manuscript,screen,review]{acmart}

\AtBeginDocument{%
  \providecommand\BibTeX{{Bib\TeX}}}

\usepackage{booktabs}
\usepackage{amsmath}

\begin{document}

%% =================== TITLE & AUTHOR BLOCK =========================
\title{AI-based Village Pond Planning System}
\subtitle{CSD Assignment 1 -- Final Technical Report}

\author{Jiya Mehta}
\affiliation{%
  \institution{IIT Bhilai / Department of Data Science and AI}
  \city{Surat}
  \country{India}}
\email{jiyam@iitbhilai.ac.in}

\renewcommand{\shortauthors}{Mehta}

%% =================== ABSTRACT ======================================
\begin{abstract}
Rural water conservation in agrarian regions relies heavily on farm ponds (\textit{talabs}) for rainwater harvesting, aquifer recharge, and drought resilience. However, conventional manual site selection and arbitrary circular buffers neglect terrain drainage, risking dry catchments or embankment washouts. This report presents an automated, distributed Web GIS system for optimal village pond siting, catchment delineation, and rainwater yield estimation. The system supports three workflows: (1) survey contour map ingestion (.kml/.kmz) interpolated via Inverse Distance Weighting (IDW), (2) bounding box selection querying a Global Digital Elevation Model (DEM) with bicubic spline zoom and contour vectorization, and (3) village name geocoding via OpenStreetMap Nominatim. Hydrological analysis runs on an 8-neighbour spherical terrain graph executing D8 steepest flow routing, flow accumulation, river/ridge safety clearance buffering, multi-factor suitability ranking, and reverse Breadth-First Search (BFS) catchment delineation. Storage volume is derived using the Rational Runoff Equation ($V = C \times R \times A$) with live Open-Meteo and NASA POWER climatology. The backend is deployed as a 4-node distributed cluster (Gateway on System 1, compute workers on Systems 2, 3, and 4) with an SQLite WAL 24-hour cache utilizing SHA-256 hashing and $0.002^\circ$ spatial quantization, reducing query latency from 3.7~s to under 90~ms.
\end{abstract}

\keywords{Village Pond Planning, Geospatial Analysis, Rainwater Harvesting, D8 Flow Routing, Catchment Delineation, Distributed Systems, Load Balancing, Web GIS}

\maketitle

%% =====================================================================
\section{Introduction}
\label{sec:intro}
Water scarcity severely restricts agricultural productivity in semi-arid and monsoonal rural regions. Torrential monsoon downpours generate rapid surface runoff and soil erosion, followed by protracted seasonal droughts. Decentralized earthen farm ponds (\textit{talabs}) capture surface runoff and replenish local water tables. However, constructing ponds without quantitative topographic analysis frequently places them in dry zones or directly inside destructive river floodways. This project presents an automated Web GIS platform that models digital elevation surfaces, ranks safe pond locations, delineates upstream catchments via flow routing, and computes expected rainwater yield using live agro-meteorological data.

\subsection{Motivation}
Manual site selection suffers from four major challenges:
\begin{itemize}
    \item \textbf{Topographic Complexity:} Visual inspection cannot reliably trace subtle ridge divides and drainage accumulation across undulating rural landscapes.
    \item \textbf{Disjoint Meteorological Data:} Local planning rarely integrates historical precipitation records or soil runoff coefficients into storage sizing.
    \item \textbf{Arbitrary Buffer Heuristics:} Conventional GIS tools often approximate drainage with circular buffers, severely miscalculating actual catchment areas.
    \item \textbf{High Survey Overhead:} On-site land surveying requires specialized tools and surveyors, hindering rapid multi-village assessment.
\end{itemize}

\subsection{Scope of the Project}
The system covers: vector contour parsing (.kml/.kmz), village geocoding, Global DEM querying, bicubic spline surface generation, D8 flow accumulation, stream safety buffering, multi-criteria pond ranking, reverse-flow BFS watershed delineation, live rainfall fetching, Rational runoff volume calculation, and distributed load balancing across 4 nodes with persistent 24h caching. It excludes civil geotechnical soil-bearing testing, structural spillway masonry design, and subterranean bedrock permeability modeling.

%% =====================================================================
\section{Problem Statement and Requirements}
\label{sec:requirements}

The system satisfies all functional requirements through modular backend and frontend components mapped in Table~\ref{tab:requirements}.

\begin{table}[h]
  \caption{Functional requirements and system implementation mapping}
  \label{tab:requirements}
  \small
  \begin{tabular}{p{3.8cm}p{3.5cm}p{6.0cm}}
    \toprule
    Requirement & Implemented in Module & Technical Implementation Details \\
    \midrule
    Base map display & \texttt{frontend/app.js} & Leaflet.js map with OSM tiles, custom markers, and layer controls. \\
    Contour map visualization & \texttt{dem\_service.py} & Matplotlib QuadContour vectorization producing dynamic GeoJSON polylines. \\
    Available land identification & \texttt{main.py}, \texttt{load\_balancer.py} & Leaflet rectangle drag bounding box or Nominatim administrative geocoding. \\
    Terrain surface generation & \texttt{terrain.py}, \texttt{dem\_service.py} & SciPy cKDTree IDW interpolation and bicubic spline grid zooming. \\
    Channel \& ridge exclusion & \texttt{catchment.py} & Exact Euclidean Distance Transform safety buffering ($D_{\text{safe}} = 200\text{ m}$). \\
    Catchment area estimation & \texttt{catchment.py} & Uphill reverse-flow Breadth-First Search (BFS) over D8 directed flow graph. \\
    Historical rainfall query & \texttt{rainfall\_service.py} & Open-Meteo Archive API with automatic fallback to NASA POWER 30-year data. \\
    Runoff volume estimation & \texttt{catchment.py} & Hydrological Rational Equation ($V = C \cdot R \cdot A$) with selectable soil catalog. \\
    Pond depth / storage recommendation & \texttt{catchment.py}, \texttt{main.py} & Multi-candidate spatial separation ranking ($m^3$, Liters, and Megaliters). \\
    Distributed load balancing & \texttt{load\_balancer.py} & Round-Robin reverse proxy with active health checks and worker failover. \\
    Multi-level caching & \texttt{cache\_service.py} & SQLite WAL 24h cache with SHA-256 file hashing and $0.002^\circ$ quantization. \\
    Combined overlay view & \texttt{frontend/app.js} & Synchronized layers (ponds, catchments, channels, contours) and metrics. \\
    \bottomrule
  \end{tabular}
\end{table}

\subsection{Non-Functional Requirements}
\begin{enumerate}
    \item \textbf{Latency:} Fresh computations must complete within 5~seconds; cached queries must return in under 100~ms.
    \item \textbf{Scalability:} Stateless compute workers must scale horizontally across multiple hosts without gateway contention.
    \item \textbf{Fault Tolerance:} The gateway must detect worker failures within 2~seconds and re-route traffic without dropping requests.
    \item \textbf{API Resilience:} Meteorological and geocoding failures must degrade gracefully via secondary endpoints.
    \item \textbf{Zero-Install Usability:} Fully accessible via any modern browser without desktop GIS software or plugins.
\end{enumerate}

%% =====================================================================
\section{System Architecture and High-Level Design}
\label{sec:architecture}

The system is deployed as a 4-node distributed cluster separating client routing, caching, and scientific computing.

Traffic enters through System 1 (\texttt{stu46\_sys1}, public port 5261, host port 5000), which operates as the Gateway and Cache Master. System 1 serves the frontend Web GIS, runs background health checks against compute nodes, executes round-robin scheduling, and manages a shared SQLite WAL 24-hour cache. 

On receiving a request, System 1 computes a normalized cache key. If found, the cached JSON payload is returned in under 90~ms without contacting external APIs or compute workers. On a cache miss, the gateway asynchronously dispatches the payload across private subnet \texttt{172.17.0.0/16} to one of three compute workers: System 2 (\texttt{172.17.0.63:5000}), System 3 (\texttt{172.17.0.64:5000}), or System 4 (\texttt{172.17.0.65:5000}). Workers execute terrain modeling, D8 flow routing, and catchment delineation, returning results to System 1 for caching and client delivery.

\subsection{Technology Stack}
\begin{itemize}
    \item \textbf{Runtime \& Framework:} Python 3.12/3.14, FastAPI with Uvicorn ASGI asynchronous server.
    \item \textbf{Scientific Computing:} NumPy (matrix operations), SciPy (cKDTree, distance transform, spline zoom), Shapely, and Matplotlib.
    \item \textbf{Frontend Web GIS:} HTML5, CSS design tokens, and Leaflet.js 1.9.4.
    \item \textbf{Persistence \& Caching:} SQLite in WAL mode with SHA-256 and spatial quantization.
    \item \textbf{Geospatial APIs:} Open-Elevation (DEM), OpenStreetMap Nominatim (Geocoding), Open-Meteo and NASA POWER (Precipitation).
\end{itemize}

%% =====================================================================
\section{Methodology}
\label{sec:methodology}

\subsection{Terrain and Elevation Surface Construction}
\subsubsection{Workflow 1: IDW from Contour Vectors}
From uploaded KML/KMZ files, contour coordinates and elevations $z_i$ are parsed. An $N \times N$ elevation surface is reconstructed using Inverse Distance Weighting (IDW):
\begin{equation}
Z(x) = \frac{\sum_{i=1}^{k} z_i / d(x, x_i)^p}{\sum_{i=1}^{k} 1 / d(x, x_i)^p}
\end{equation}
where $d(x, x_i)$ is the Euclidean distance queried via a SciPy \texttt{cKDTree} ($k=12$, power $p=2$).

\subsubsection{Workflows 2 and 3: Global DEM Fetching and Spline Interpolation}
For bounding boxes and village search, elevations are queried from the Global DEM API and upsampled to a dense $50 \times 50$ matrix using third-order bicubic spline interpolation (\texttt{scipy.ndimage.zoom}):
\begin{equation}
Z(r, c) = \sum_{i=0}^{3} \sum_{j=0}^{3} a_{ij} (r - r_0)^i (c - c_0)^j
\end{equation}
Elevation contours are extracted from this matrix via marching squares (\texttt{matplotlib.pyplot.contour}).

\subsection{Catchment Area Delineation}
Drainage follows the path of steepest downhill descent. On the elevation grid, an 8-neighbour graph is constructed. The distance $d_{\text{hav}}(u, v)$ between cells is computed via the Haversine formula to eliminate latitudinal distortion:
\begin{equation}
d_{\text{hav}}(u, v) = 2 R_{\text{earth}} \arcsin \left( \sqrt{\sin^2(\Delta \phi / 2) + \cos \phi_u \cos \phi_v \sin^2(\Delta \lambda / 2)} \right)
\end{equation}
Slope gradient $S(u, v) = (Z(u) - Z(v)) / d_{\text{hav}}(u, v)$ determines the downhill flow edge:
\begin{equation}
v^* = \arg\max_{v \in \mathcal{N}_8(u)} S(u, v) \quad \text{for } S(u, v^*) > 0
\end{equation}
Cells with $S \le 0$ for all neighbors are identified as local terrain depressions or sinks. Flow accumulation $A_{\text{acc}}(u)$ is computed in $\mathcal{O}(|\mathcal{V}| + |\mathcal{E}|)$ by traversing nodes in reverse topological order:
\begin{equation}
A_{\text{acc}}(u) = 1 + \sum_{v \in \text{InEdges}(u)} A_{\text{acc}}(v)
\end{equation}

Cells in the 98th percentile of accumulation define major drainage channels. To prevent embankment washouts, candidates within safety distance $D_{\text{safe}} = 200\text{ m}$ of streams are excluded using an Exact Euclidean Distance Transform:
\begin{equation}
D_{\text{channel}}(x) = \min_{y \in \mathcal{V}_{\text{river}}} \|x - y\|_2
\end{equation}
Safe cells are ranked via multi-criteria suitability scoring:
\begin{equation}
\text{Score}(u) = 45 \cdot \frac{\ln(1 + A_{\text{acc}}(u))}{\ln(1 + A_{\max})} + 15 \cdot \frac{Z_{\max} - Z(u)}{Z_{\max} - Z_{\min}} + 35 \cdot \min\left(1.0, \frac{\bar{Z}_{\text{adj}}(u) - Z(u)}{\Delta Z_{\text{relief}}}\right)
\end{equation}
Spatial separation $\|u - v\|_2 \ge \delta_{\text{sep}}$ ensures distinct candidate sites. Contributing catchments are delineated via reverse BFS from each selected sink:
\begin{equation}
\mathcal{C}(u_{\text{pond}}) = \{u_{\text{pond}}\} \cup \{v \in \mathcal{V} \mid \text{Path}(v \rightsquigarrow u_{\text{pond}}) \text{ in } \mathcal{G}\}
\end{equation}
Total catchment area $A$ is obtained by summing spherical cell surface areas within $\mathcal{C}$.

\subsection{Rainfall Integration and Rational Runoff Estimation}
Annual precipitation $R$ (in meters) is queried from Open-Meteo (with fallback to NASA POWER). Runoff volume $V$ is estimated via the Rational Equation:
\begin{equation}
V = C \times R \times A
\end{equation}
where $C$ is the soil runoff coefficient ($0.30$ for vegetative sandy loam, $0.40$ for clay loam agricultural soil, $0.50$ for hard clay, and $0.60$ for barren clay).

%% =====================================================================
\section{Implementation}
\label{sec:implementation}

\subsection{Backend and API Design}
The gateway exposes modular REST endpoints documented in Table~\ref{tab:api}.

\begin{table}[h]
  \caption{Gateway REST API Endpoints}
  \label{tab:api}
  \small
  \begin{tabular}{llp{7.5cm}}
    \toprule
    HTTP Method & Endpoint Path & Operational Purpose \\
    \midrule
    \texttt{GET} & \texttt{/} & Serves the Leaflet GIS Web Application. \\
    \texttt{GET} & \texttt{/health} & Gateway health check and active worker count. \\
    \texttt{GET} & \texttt{/cluster/status} & Real-time worker cluster latencies, request counts, and cache metrics. \\
    \texttt{GET} & \texttt{/api/rainfall} & Fetches precipitation for coordinates with 24h spatial caching. \\
    \texttt{GET} & \texttt{/api/geocode} & Geocodes place name to bounding box coordinates via Nominatim. \\
    \texttt{POST} & \texttt{/analyzeContour} & Case 1: Ingests multipart .kml/.kmz file, checks SHA-256 cache, routes to worker. \\
    \texttt{POST} & \texttt{/analyzeArea} & Case 2: Ingests bounding box, quantizes coordinates, routes DEM analysis. \\
    \texttt{POST} & \texttt{/analyzePlace} & Case 3: Ingests village name, geocodes, normalizes string, routes DEM analysis. \\
    \bottomrule
  \end{tabular}
\end{table}

\subsection{Frontend and Visualization}
The web interface is built using Leaflet.js and responsive CSS. It provides tabbed selection across the three workflows, quick village chips, interactive bounding box drawing, and crosshair coordinate tracking. Visual layers display ranked pond markers ($P_1, P_2, \dots$), upstream catchment boundary polygons, natural stream polylines, elevation contour vectors, and an interactive dashboard showing coordinates, elevation, catchment area (ha), and expected water volume ($m^3$, Liters, Megaliters).

\subsection{Database and Storage}
The cache is implemented in \texttt{cache\_service.py} using an embedded SQLite database configured in Write-Ahead Logging (\texttt{PRAGMA journal\_mode=WAL}) with B-tree indexes on \texttt{cache\_key} and \texttt{expires\_at}. Caching logic is tailored to each workflow:
\begin{itemize}
    \item \textbf{Case 1 (Files):} Input bytes are hashed using SHA-256 with analysis parameters.
    \item \textbf{Case 2 (Bounding Box):} Coordinates are quantized to $0.002^\circ$ ($\approx 200\text{ m}$), mapping minor dragging variations to identical cache keys.
    \item \textbf{Case 3 (Places):} Village names are sanitized, lowercased, and converted to uniform slugs.
\end{itemize}

%% =====================================================================
\section{CSD Themes and Topics Applied in the Project}
\label{sec:csd-themes}

The system architecture embodies core Computer System Design principles mapped in Table~\ref{tab:csd-themes}.

\begin{table*}[t]
  \caption{Mapping of core CSD themes/topics to their use in this project}
  \label{tab:csd-themes}
  \small
  \begin{tabular}{p{3.2cm}p{4.6cm}p{7.4cm}}
    \toprule
    CSD Theme/Topic & Where used in the project & Justification / design rationale \\
    \midrule
    API Design (REST) & \texttt{main.py}, \texttt{load\_balancer.py} & Standardized HTTP verbs, explicit error codes (200, 400, 502), typed Pydantic payloads, and clean OpenAPI decoupling between UI and workers. \\
    Load Balancing & \texttt{load\_balancer.py} on System 1 & Distributes CPU-intensive terrain calculations across Systems 2, 3, and 4 via Round-Robin with least-connections bias to prevent worker saturation. \\
    Caching & \texttt{cache\_service.py} on System 1 & SQLite WAL persistent 24h cache eliminates redundant external DEM/meteo queries and expensive IDW calculations ($3.7\text{s} \to <90\text{ms}$). \\
    Database Indexing / Query Optimization & \texttt{cache\_service.py} & B-tree indexes on \texttt{cache\_key} and \texttt{expires\_at} guarantee $\mathcal{O}(1)$ query lookups and rapid pruning of expired records. \\
    Concurrency / Asynchronous Processing & \texttt{load\_balancer.py} & Python AsyncIO with \texttt{httpx.AsyncClient} prevents blocking the gateway event loop while awaiting remote worker responses. \\
    Microservices vs. Monolith & Gateway/Worker Split & Decoupled master-worker design: System 1 gateway remains ultra-responsive for static UI serving while stateless worker nodes scale horizontally. \\
    Design Patterns & Reverse Proxy, Strategy, Factory & Reverse proxy handles traffic dispatching; Strategy pattern enables Open-Meteo to NASA POWER fallback; Factory instantiates terrain generators. \\
    Authentication / Authorization & CORS \& IP Whitelisting & Internal network whitelisting restricts worker execution to System 1 gateway requests, preventing unauthorized external access. \\
    Containerization / Deployment & Python venv \& Daemons & Isolated environments on each lab host with background execution (\texttt{nohup}) ensure reproducible library dependencies. \\
    Error Handling and Resilience & In-flight failover \& retries & Gateway detects connection drops mid-flight and transparently re-routes requests to healthy workers, preventing user-facing 502 errors. \\
    Algorithms and Complexity & D8 Flow \& Reverse BFS & D8 topological sort ($\mathcal{O}(V \log V)$) and reverse BFS ($\mathcal{O}(V + E)$) chosen over naive $\mathcal{O}(N^2)$ all-pairs traversal for sub-second execution. \\
    Testing Strategy & Pytest Unit \& Cluster Tests & 15 automated test cases in \texttt{tests/} validate geometry, parsing, flow algorithms, and end-to-end cluster failover. \\
    Version Control / CI-CD & Git Repository & Structured Git tracking on GitHub coordinated deployment across all 4 server nodes. \\
    \bottomrule
  \end{tabular}
\end{table*}

\paragraph{Significant Engineering Decisions}
Two architectural choices were critical to system efficacy:
\begin{enumerate}
    \item \textbf{Decoupled Master-Worker vs. Monolith:} A monolithic server executing dense matrix IDW calculations suffered severe socket saturation under simultaneous queries. Offloading compute to Systems 2, 3, and 4 keeps the System 1 Gateway CPU near zero, ensuring consistent sub-millisecond static page serving and proxy routing.
    \item \textbf{Intelligent Quantized Caching vs. Exact Keys:} Exact coordinate matching fails in interactive Web GIS because human mouse dragging rarely produces identical coordinates. Quantizing bounding coordinates to $0.002^\circ$ transformed the cache into a practical system delivering $<100$~ms responses on repeated queries.
\end{enumerate}

%% =====================================================================
\section{Results and Evaluation}
\label{sec:results}

The cluster was validated through end-to-end operational testing across five scenarios on the university network (\texttt{10.1.75.79:5261}), summarized in Table~\ref{tab:results}.

\begin{table}[h]
  \caption{Representative Hydrological Results Across Tested Locations}
  \label{tab:results}
  \small
  \begin{tabular}{lcccccc}
    \toprule
    Test Case / Location & Lat, Lon & Rain ($R$) & Runoff ($C$) & Catchment ($A$) & Top Water Vol ($V$) & Execution Mode \\
    \midrule
    Case 1: \texttt{contours\_1m.kml} & $21.2528^\circ, 81.2872^\circ$ & 1289 mm & 0.40 & 1.631 ha & $8,410.5\text{ m}^3$ (8.41 ML) & Worker Compute \\
    Case 2: Selected Farm Area & $21.2486^\circ, 81.3006^\circ$ & 1289 mm & 0.40 & 9.407 ha & $48,514.6\text{ m}^3$ (48.51 ML) & Worker Compute \\
    Case 3: Khapri Village & $21.0533^\circ, 79.0451^\circ$ & 1052 mm & 0.40 & 1.778 ha & $7,481.9\text{ m}^3$ (7.48 ML) & Worker Compute \\
    Case 3: Khapri (Re-query) & $21.0533^\circ, 79.0451^\circ$ & 1052 mm & 0.40 & 1.778 ha & $7,481.9\text{ m}^3$ (7.48 ML) & 24h Cache Hit \\
    Case 3: Bhilai Urban Fringe & $21.2167^\circ, 81.4333^\circ$ & 1289 mm & 0.40 & 40.710 ha & $209,917.0\text{ m}^3$ (209.9 ML) & Worker Compute \\
    \bottomrule
  \end{tabular}
\end{table}

\subsection{Performance}
Latency benchmarks demonstrated high throughput under cold compute versus warm cache conditions:
\begin{itemize}
    \item \textbf{Case 1 (KML Upload):} Cold compute completed in 2,880~ms; identical file re-upload returned via SHA-256 cache in \textbf{216~ms}.
    \item \textbf{Case 2 (Area DEM):} Fresh DEM spline interpolation took 2,769~ms; re-selection within quantized bounds returned in \textbf{180~ms}.
    \item \textbf{Case 3 (Village Search):} Fresh query for Khapri took 3,867~ms on Worker 1 (\texttt{sys2}); subsequent query responded from System 1 cache in \textbf{89.5~ms} (a $43\times$ speedup).
    \item \textbf{Round-Robin Load Distribution:} Verification confirmed sequential request distribution across workers: Request 1 $\to$ \texttt{sys2}, Request 2 $\to$ \texttt{sys3}, Request 3 $\to$ \texttt{sys4}.
\end{itemize}

%% =====================================================================
\section{Discussion and Limitations}
\label{sec:discussion}
While the system successfully addresses the requirements, three engineering limitations are noted:
\begin{enumerate}
    \item \textbf{DEM Spatial Resolution:} Public Open-Elevation DEM data operates at $\approx 30$-meter resolution. While bicubic spline interpolation smooths elevations, subtle micro-topographic field bunds ($<1$~meter height) are not captured unless surveyed drone LIDAR data is supplied via Case 1 KML.
    \item \textbf{Constant Runoff Coefficients:} The Rational equation assumes constant runoff coefficients throughout a rain event. Integrating dynamic Green-Ampt infiltration or SCS Curve Number methods would better account for antecedent soil moisture across multi-day monsoonal storms.
    \item \textbf{Gateway Redundancy:} System 1 acts as a single point of network entry for port 5261. In enterprise production, a Virtual IP (VIP) managed via Keepalived or an AWS Elastic Load Balancer would provide redundant multi-master failover.
\end{enumerate}

%% =====================================================================
\section{AI Tool Usage Declaration}
\label{sec:ai-usage}
In compliance with the assignment's LLM Usage Policy, generative AI tools (including ChatGPT and Claude) were utilized as technical sounding boards and productivity accelerators during specific stages of project development:
\begin{itemize}
    \item \textbf{Conceptual Discussions and Algorithmic Logic:} AI assistance was used for discussing mathematical formulations, specifically reviewing D8 steepest-gradient directional logic, verifying spherical Haversine trigonometric distance equations, and exploring standard soil-runoff coefficient ranges for the Rational Method.
    \item \textbf{Code Improvement and Debugging:} AI tools were consulted to identify and resolve complex syntax and runtime bugs, such as multipart form-data parameter mismatches between client and server, non-blocking asynchronous socket timeouts in \texttt{httpx}, and database concurrency locks under multi-threaded SQLite execution.
    \item \textbf{Frontend UI Assistance:} AI was utilized to suggest modern CSS layout structures, responsive flexbox card arrangements, and Leaflet layer group management patterns.
    \item \textbf{Report Formatting:} AI was used to format technical content into the standard ACM \texttt{acmart} manuscript LaTeX template.
    \item \textbf{Manual Verification and Testing:} All core system architecture, distributed cluster configurations across the four lab nodes, mathematical pipelines, and algorithmic implementations were manually reviewed, validated, and thoroughly tested using automated pytest suites and end-to-end browser verification by the student.
\end{itemize}

%% =====================================================================
\appendix

\section{Source Code and Repository}
The complete source code, automated test suite, frontend assets, and deployment scripts are available at:
\begin{center}
\url{https://github.com/jiyaa25/csd_asg_pond}
\end{center}

\subsection{Repository Folder Structure}
\begin{itemize}
    \item \texttt{load\_balancer.py}: System 1 Gateway, health daemon, worker load balancer, and caching router.
    \item \texttt{cache\_service.py}: SQLite Write-Ahead Logging (WAL) cache engine with spatial quantization.
    \item \texttt{main.py}: Core FastAPI worker computation service running on Systems 2, 3, and 4.
    \item \texttt{catchment.py}: D8 flow accumulation, river/ridge exclusion, and uphill reverse-flow BFS.
    \item \texttt{terrain.py} \& \texttt{graph.py}: IDW terrain interpolation and spherical 8-neighbour surface graphs.
    \item \texttt{dem\_service.py}: Global DEM querying, bicubic spline zooming, and contour extraction.
    \item \texttt{rainfall\_service.py}: Dual-fallback precipitation service (Open-Meteo and NASA POWER).
    \item \texttt{frontend/}: Interactive Leaflet.js Web GIS application (\texttt{index.html}, \texttt{app.js}, \texttt{style.css}).
    \item \texttt{tests/}: Automated pytest unit and integration test suite (\texttt{test\_all.py}, \texttt{test\_parser.py}).
    \item \texttt{report.tex}: Complete ACM manuscript LaTeX source code for this technical report.
\end{itemize}

\end{document}
\endinput
